\label{sec:conclude}

In this paper, we introduced \system, a database-backed mistake-memory framework for LLM-based CWE vulnerability detection. The main problem we addressed is that standard LLM inference treats each code sample independently and does not preserve previous failures. As a result, the same model can repeatedly make similar false-positive or false-negative predictions. To address this issue, \system stores incorrect predictions in a persistent SQLite database and retrieves similar past mistakes during future inference. The retrieved records are inserted into the prompt as cautionary examples, allowing the model to consider its previous failures without fine-tuning or changing the base model architecture.

Our evaluation used Qwen2.5-Coder-7B-Instruct on CASTLE-C250 and a sampled Juliet C/C++ Test Suite subset. CASTLE provides a controlled benchmark with 250 hand-crafted C micro-benchmarks covering 25 common CWE categories, while Juliet provides a larger external benchmark organized under 118 CWE categories~\cite{castle,castle_github,juliet_nist}. The results show that \system has mixed behavior on CASTLE. On the held-out CASTLE split, accuracy and F1 decreased because false negatives increased. On the full CASTLE in-dataset analysis, \system reduced false positives from 53 to 31, but increased false negatives from 30 to 61. This shows that the original prompt can make the model more conservative. However, on the external Juliet subset, the CASTLE-built mistake database improved accuracy from 68.67\% to 82.40\% and macro F1 from 66.23\% to 82.18\%. We also added a supplemental full-C250 compact-RAG experiment with OpenRouter's \texttt{openai/gpt-oss-120b:free}. In that setting, \system improved held-out F1 from 0.616 to 0.709 and CWE correctness from 0.433 to 0.520. This suggests that persistent mistake memory can transfer useful information across CWE benchmarks and can improve the original benchmark when retrieval is selective enough.

\noindent\textbf{Limitations and Future Work.}
The current prototype has several limitations. First, most of the original experiments use one base model because of runtime constraints, and the OpenRouter experiment uses a free API endpoint with unstable latency. Testing more LLMs under the same full-C250 protocol would help determine whether \system is model-specific or generally useful across different code models. Second, the current retrieval method uses TF-IDF similarity and lightweight RAG reranking, which mainly capture lexical overlap. Future work could use code embeddings or a vector database to retrieve semantically similar mistakes more accurately. Third, the current prompt does not carefully balance false-positive and false-negative examples. This caused \system to reduce over-reporting in some settings but also increased under-reporting on CASTLE. A better design could retrieve separate sets of false-positive and false-negative records, use compact RAG evidence packs, or assign different weights to each type of mistake. Finally, our evaluation focuses on CWE vulnerability detection. The same architecture could be extended to other LLM-based software engineering tasks, such as program repair, code review, or test generation, by changing the mistake schema and retrieval representation.
